\underline{Net Present Value}

\begin{definition}
    Future payments from an initial investment are called \emph{cash flows}.
\end{definition}

\begin{definition}
    \emph{Net present value (NPV)} is the sum of the present value of the cash flows minus the initial investment.
\end{definition}

NOTE: \(NPV>0\) means the investment is ultimately profitable while \(NPV<0\) means that the investment is a loss.

\begin{example}
    An initial investment of \(\$1000\) guarantees cash flows following the table. If an interest of \(3\%\) compounding annually is applied, find the net present value.
    \begin{table}[]
        \centering
        \begin{tabular}{c|c}
            year & cash flow \\\hline
             \(1\)&\(500\) \\
             \(2\)& \(450\) \\
             \(3\)& \(400\)
        \end{tabular}
        \caption{Payments at the end of each of the first \(3\) years}
    \end{table}

    \begin{solution}
        \[NPV = 500(1.03)^{-1}+450(1.03)^{-2}+400(1.03)^{-3}  - 1000 \approx 275.66.\]
        A profit. That is, you would have to invest \(500(1.03)^{-1}+450(1.03)^{-2}+400(1.03)^{-3} = 1275.66\) in order to get the same return that you got with the \(\$1000\) investment.
    \end{solution}
\end{example}